\documentclass[12pt]{article}

% Packages
\usepackage[margin=1in]{geometry}
\usepackage{amsmath}
\usepackage{booktabs}
\usepackage{array}
\usepackage{hyperref}
\usepackage{natbib}
\usepackage{setspace}
\usepackage{titlesec}
\usepackage{parskip}
\usepackage[T1]{fontenc}
\usepackage[utf8]{inputenc}
\usepackage{lmodern}

\bibpunct{(}{)}{;}{a}{,}{,}
\setstretch{1.5}

% Title
\title{\textbf{Relative Performance Feedback and Redistribution Preferences: Experimental Evidence}}
\author{Isabella Valencia\\[6pt]
  \small Paris School of Economics\\[2pt]
  \small \textit{Supervised by Prof.\ Andrea FM Martinangeli}}
\date{}

\begin{document}

\maketitle

\begin{abstract}
This paper examines how relative performance information shapes individuals' redistribution decisions. Using an online experiment, participants first complete a real-effort task and are randomly assigned to one of two feedback conditions: no information about their relative standing, or knowledge of whether they performed above or below a randomly paired peer. They then decide how to allocate a fixed sum between two workers whose earnings were determined either by merit (task performance) or by luck (lottery). We elicit beliefs about the role of effort versus luck in income generation, allowing us to study perceptions as a potential mechanism. The design allows us to (i) isolate the causal effect of relative position feedback on redistribution, (ii) compare redistribution under merit versus luck framing within subjects, and (iii) link self-serving attribution patterns to allocation decisions.
\end{abstract}

\newpage

%-----------------------------------------------------------------------
\section{Introduction}
%-----------------------------------------------------------------------

How much individuals are willing to redistribute from those who earn more to those who earn less is one of the central questions in political economy. A large body of evidence documents that support for redistribution is not driven solely by material self-interest: fairness considerations, and in particular beliefs about whether inequality reflects merit or chance, play a decisive role \citep{alesina2005, fong2001, benabou2006}. When people believe that effort and ability determine economic outcomes, they tend to accept larger inequalities. When they attribute outcomes to luck or circumstance, they support more redistribution \citep{piketty1995, cappelen2007}.

Yet the social context in which individuals form these beliefs has received comparatively little attention. In reality, people rarely evaluate inequality in a vacuum: they are constantly exposed to information about their own relative standing — in earnings, education, or professional achievement. The question of whether this relative position information affects redistribution preferences, and through what mechanism, remains largely open.

This paper addresses that gap. We hypothesize that receiving information about one's relative performance generates motivated reasoning \citep{kunda1990}: individuals who learn that they outperformed a peer are more likely to attribute the inequality to merit, while those who learn they underperformed are more likely to attribute it to luck or circumstance. Both responses are self-serving — the ``winner'' justifies keeping more, and the ``loser'' justifies taking more — but they operate through the same underlying channel: perceived fairness of the income-generating process. Evidence from our experiment suggests that this asymmetry may be stronger for those at the top of the comparison than for those at the bottom.

We test this in a pre-registered online experiment with three components. In \textbf{Part~1}, participants complete a code recognition task (real-effort, no prior knowledge required) across five timed rounds. Before seeing any results, they answer incentive-compatible questions about their effort, their expected score, and their expected relative position. They are then randomly assigned to one of two feedback conditions: a \textbf{no-information group} that receives only a completion message, or a \textbf{relative comparison group} that learns whether they performed above or below a randomly drawn peer. In \textbf{Part~2}, participants make two sequential allocation decisions. In the first, they distribute \$8 between two workers whose earnings were determined by task performance (merit framing). In the second — presented as a separate case — they distribute \$8 between two workers whose earnings were determined by lottery (luck framing). This within-subjects design isolates the effect of income source on redistribution while holding individual heterogeneity constant. Finally, participants answer a question about the extent to which they believe income in the task was determined by effort versus luck, which we use as a mediating variable.

Our design contributes to three strands of literature. First, we add to the experimental economics literature on fairness ideals and redistribution, building directly on the framework developed by \citet{cappelen2007} and \citet{almas2020}. \citet{cappelen2007} introduced the spectator paradigm for studying distributive preferences: participants act as impartial third-party allocators who observe income generated under different processes and decide how to redistribute it. Using this approach, they documented systematic heterogeneity in fairness views — egalitarian, meritocratic, and libertarian — and showed that the income-generating process (merit vs.\ luck) strongly affects redistribution decisions. \citet{almas2020} extended this framework to cross-national comparison, using a real-effort code recognition task and a between-subjects design in which American and Norwegian spectators were randomly assigned to observe inequality arising from either a merit or a luck condition. They found that meritocratic beliefs are prevalent in both countries but stronger in the United States, and that the merit vs.\ luck manipulation produces large differences in the inequality spectators are willing to accept.

Our experiment is built directly on this design. We adopt the same code recognition task introduced in \citet{almas2020}, which requires no domain knowledge and ensures variation in performance reflects effort and attention. We also adopt their discrete-choice allocation format \citep{cappelen2007}, presenting participants with a menu of options in \$1 increments between full preservation of inequality and full equalization. Our primary focus, like theirs, is on the merit condition — the case in which income differences reflect real performance. Where we depart from their design is in (i) adding a between-subjects social comparison treatment orthogonal to the income-generating process, and (ii) using a within-subjects design in which the merit decision precedes a luck decision. The luck decision is included not as an independent treatment for drawing causal inferences about luck-vs-merit differences, but as a within-person benchmark: by observing each participant's redistribution in both the merit and luck frames, we can characterize individual heterogeneity in how much more a person is willing to redistribute when income is luck-determined, and use this as a control for their underlying fairness type. The merit case remains the primary outcome of interest.

Second, we contribute to the literature on social comparisons and economic behavior \citep{clark1996, fliessbach2007, card2012}. Much of this work focuses on the effect of relative income on individual well-being or effort provision; we instead focus on its downstream effect on distributional preferences. Third, we connect to the motivated reasoning and self-serving bias literature in psychology and behavioral economics \citep{konow2000, dahl1999, haisley2010}. Knowing one is ``above'' versus ``below'' a peer may shift how one attributes the peer's income — and by extension, how one evaluates the fairness of inequality.

The remainder of the paper is organized as follows. Section~\ref{sec:design} describes the experimental design. Section~\ref{sec:hypotheses} presents the theoretical predictions and hypotheses. Section~\ref{sec:empirics} discusses the empirical strategy. Section~\ref{sec:conclusion} concludes.

%-----------------------------------------------------------------------
\section{Experimental Design}
\label{sec:design}
%-----------------------------------------------------------------------

\subsection{Overview and Sample}

The experiment is conducted online using JATOS and takes approximately 7--10 minutes to complete. Participants are recruited through [platform]. The study is available in English and Spanish to broaden geographic reach. All data are collected anonymously under GDPR regulations. Participation is voluntary and non-compensated.

\subsection{Part 1: Real-Effort Task}

Participants complete a \textbf{code recognition task} adapted from \citet{almas2020}. The task is designed to require no prior knowledge, ensuring that variation in performance reflects effort and attention rather than domain-specific skills. In each round, a three-digit target code is displayed and participants must identify and click all instances of that code within a matrix of three-digit numbers. The matrix contains approximately 320 cells (20$\times$16 on desktop, 16$\times$20 on mobile).

The task consists of \textbf{five rounds of 45 seconds each} (total: 3 minutes 45 seconds), with a different target code in each round. Before the real task, participants complete a \textbf{20-second practice round} (not counted toward their score) to familiarize themselves with the interface.

\textbf{Scoring}: $+1$ point for each correct marking, $-1$ point for each incorrect marking, and no penalty for missing a target. The theoretical maximum score is approximately 200 points. This incentive structure penalizes guessing and rewards precise performance.

\subsection{Pre-Feedback Self-Report}

Immediately after the task, and \textbf{before receiving any performance feedback}, participants answer three questions designed to elicit their beliefs and effort perceptions:

\begin{enumerate}
  \item \textbf{Effort} (slider, 0--10): ``How much effort did you put into the task?''
  \item \textbf{Score estimate} (integer, $-200$ to 200): ``The maximum amount of points you could have obtained was 200. How many do you think you got?''
  \item \textbf{Relative performance} (binary): ``If you were randomly paired with another participant, do you think you would likely perform better or worse than them?'' (\textit{better or the same} / \textit{worse})
\end{enumerate}

These questions are posed before feedback to avoid contamination of beliefs by the treatment and to enable analysis of prediction errors and overconfidence.

\subsection{Feedback Treatment (Between-Subjects)}

After the self-report, participants are randomly assigned — at page load, before any interaction — to one of two feedback conditions:

\textbf{No-information group} (1/3 of participants): Participants see only a message confirming completion of Part~1 and are directed to Part~2. They receive no information about their score or relative performance.

\textbf{Relative comparison group} (2/3 of participants): Participants learn whether their score was \textbf{above or below} the score of a randomly drawn previous participant. The pair score is drawn from a fixed pool of two anchor values representing a strong performer and a weak performer. Critically, participants are shown only the directional comparison (above/below) — they do not see the pair's actual score or their own score. This design cleanly isolates the informational content of relative standing from the magnitude of the performance gap.

The asymmetry in group sizes (1/3 vs.\ 2/3) is chosen to maximize statistical power for the comparison of interest, which involves the relative feedback condition.

\subsection{Part 2: Allocation Decisions}

Part~2 consists of two sequential, incentivized allocation tasks. In each task, the participant acts as an impartial third-party allocator deciding how to distribute \$8 between two fictitious workers. Participants are told that their decision is anonymous and that one randomly selected participant's decision will be implemented as the actual payment.

\textbf{Decision 1 — Merit scenario}: Participants are told that Worker~A and Worker~B completed the same task the participant just completed. Worker~A achieved the highest score and initially received \$8; Worker~B received \$0. Participants choose how to distribute the \$8 between the two workers.

\textbf{Decision 2 — Luck scenario}: Participants are told that Worker~C and Worker~D also completed the task. In this case, initial earnings were assigned by lottery with no regard for performance: Worker~C won the lottery and received \$8; Worker~D received \$0. Participants again choose how to distribute the \$8.

The order of the two decisions is fixed: the merit scenario always precedes the luck scenario. This ordering reflects the study's primary focus: the merit decision is the main outcome of interest, and we present it first so it is evaluated on its own terms, before participants have been primed to think about luck-based inequality. The luck decision is collected second as a within-subjects reference point. Rather than serving as an independent treatment for drawing causal conclusions about the luck-vs-merit distinction, the luck condition allows us to characterize each participant's baseline redistribution under a norm-neutral (luck-only) framing, and to measure individual heterogeneity in how much more they redistribute in the luck versus merit case. This within-person gap serves as a control for underlying fairness type in the analysis of treatment effects on the merit decision.

In both decisions, participants choose from a discrete menu of \textbf{nine options} in \$1 increments, following the allocation format of \citet{cappelen2007}:

\begin{table}[h]
\centering
\caption{Allocation options presented to participants}
\label{tab:options}
\begin{tabular}{lcc}
\toprule
Option & Worker A (or C) & Worker B (or D) \\
\midrule
1 (no redistribution) & \$8 & \$0 \\
2 & \$7 & \$1 \\
3 & \$6 & \$2 \\
4 & \$5 & \$3 \\
5 & \$4 & \$4 \\
6 & \$3 & \$5 \\
7 & \$2 & \$6 \\
8 & \$1 & \$7 \\
9 & \$0 & \$8 \\
\bottomrule
\end{tabular}
\end{table}

This format makes the trade-off between equity and efficiency transparent to participants and enables direct comparison of individual redistribution choices across the two frames.

\subsection{Background Questions and Belief Elicitation}

After Decision~1 and before Decision~2, participants answer a set of background questions: education level (own and parents'), country of origin (in addition to country of residence, collected before the task), occupation status and field, and any comments.

A key additional question elicits the participant's \textbf{beliefs about the role of effort versus luck} in generating the income distribution they observed:

\begin{quote}
``To what extent do you think the initial income assigned to the participants in the previous task — before any redistribution — was determined by effort versus luck?''
\end{quote}

Participants respond on a slider from 0 (100\% luck) to 100 (100\% effort). This variable serves as a \textbf{potential mediator}: if relative performance feedback shifts redistribution decisions by updating beliefs about the role of effort in income generation, then the effect of the treatment on allocation should be attenuated or eliminated when conditioning on this belief measure.

\subsection{Experimental Timeline}

\begin{table}[h]
\centering
\caption{Experimental timeline}
\label{tab:timeline}
\begin{tabular}{cll}
\toprule
Step & Screen & Key action \\
\midrule
1  & Welcome                    & Consent, honesty statement \\
2  & Early demographics         & Gender, age, country of residence \\
3  & Instructions               & Task rules explained \\
4  & Example + Practice         & 20-second practice round \\
5  & Real task                  & 5 rounds $\times$ 45 seconds \\
6  & Self-report                & Effort, score guess, relative performance belief \\
7  & \textbf{Feedback (randomized)} & None / Relative comparison \\
8  & \textbf{Decision 1}        & Allocate \$8 — merit scenario \\
9  & Late demographics + belief & Education, occupation, effort-vs-luck slider \\
10 & \textbf{Decision 2}        & Allocate \$8 — luck scenario \\
11 & Completion                 & Results displayed; study closes after 20 seconds \\
\bottomrule
\end{tabular}
\end{table}

%-----------------------------------------------------------------------
\section{Hypotheses}
\label{sec:hypotheses}
%-----------------------------------------------------------------------

Drawing on motivated reasoning theory \citep{kunda1990} and the self-serving bias in fairness judgments \citep{konow2000, haisley2010}, we state the following hypotheses:

\textbf{H1 — Income source and redistribution}: Participants will redistribute more in the luck scenario (Decision~2) than in the merit scenario (Decision~1). This replicates the canonical finding that people apply different fairness norms depending on whether inequality is perceived as deserved \citep{cappelen2007, almas2020}.

\textbf{H2 — Relative position and redistribution}: Participants who learn they performed \textit{below} their peer will allocate more to the lower-earning worker than participants who learn they performed \textit{above} their peer. Participants with no feedback will fall between the two groups.

\textbf{H3 — Self-serving attribution}: Participants who learn they performed \textit{above} their peer will report higher effort-vs-luck scores (perceiving income as more effort-driven) than those who learn they performed \textit{below}. This would be consistent with a motivated reasoning account.

\textbf{H4 — Mediation}: The effect of relative feedback on redistribution (H2) will be partially mediated by effort-vs-luck beliefs (H3). Participants who perceive income as luck-driven will redistribute more regardless of their relative position.

%-----------------------------------------------------------------------
\section{Empirical Strategy}
\label{sec:empirics}
%-----------------------------------------------------------------------

\textit{(To be developed once data are collected.)}

The primary analysis will estimate the effect of relative feedback on allocation decisions using OLS regression with the amount allocated to the lower-earning worker as the outcome variable. The main specification will be:

\begin{equation}
  y_i = \alpha + \beta \cdot \text{RelativeFeedback}_i + \gamma X_i + \varepsilon_i
  \label{eq:main}
\end{equation}

where $\text{RelativeFeedback}_i$ is an indicator for assignment to the relative comparison group, and $X_i$ is a vector of pre-determined controls (gender, age, education, country of residence). To test H2 more precisely, the relative feedback group will be split into above/below subgroups.

The within-subjects comparison of Decisions~1 and~2 (H1) will be estimated using a paired $t$-test and a fixed-effects panel regression that absorbs individual heterogeneity.

Mediation analysis (H4) will follow the causal mediation framework of \citet{imai2010}, using effort-vs-luck beliefs as the mediating variable.

%-----------------------------------------------------------------------
\section{Conclusion}
\label{sec:conclusion}
%-----------------------------------------------------------------------

\textit{(To be developed once data are collected.)}

%-----------------------------------------------------------------------
% References
%-----------------------------------------------------------------------
\bibliographystyle{apalike}
\begin{thebibliography}{99}

\bibitem[Alesina and Angeletos(2005)]{alesina2005}
Alesina, A., \& Angeletos, G.-M. (2005).
Fairness and redistribution.
\textit{American Economic Review}, 95(4), 960--980.

\bibitem[Alm{\aa}s et al.(2020)]{almas2020}
Alm{\aa}s, I., Cappelen, A.~W., \& Tungodden, B. (2020).
Cutthroat capitalism versus cuddly socialism: Are Americans more meritocratic and efficiency-seeking than Scandinavians?
\textit{Journal of Political Economy}, 128(5), 1753--1788.

\bibitem[B{\'e}nabou and Tirole(2006)]{benabou2006}
B{\'e}nabou, R., \& Tirole, J. (2006).
Belief in a just world and redistributive politics.
\textit{Quarterly Journal of Economics}, 121(2), 699--746.

\bibitem[Cappelen et al.(2007)]{cappelen2007}
Cappelen, A.~W., Hole, A.~D., S{\o}rensen, E.~{\O}., \& Tungodden, B. (2007).
The pluralism of fairness ideals: An experimental approach.
\textit{American Economic Review}, 97(3), 818--827.

\bibitem[Cappelen et al.(2013)]{cappelen2013}
Cappelen, A.~W., S{\o}rensen, E.~{\O}., \& Tungodden, B. (2013).
When do we lie?
\textit{Journal of Economic Behavior \& Organization}, 93, 258--265.

\bibitem[Card et al.(2012)]{card2012}
Card, D., Mas, A., Moretti, E., \& Saez, E. (2012).
Inequality at work: The effect of peer salaries on job satisfaction.
\textit{American Economic Review}, 102(6), 2981--3003.

\bibitem[Clark and Oswald(1996)]{clark1996}
Clark, A.~E., \& Oswald, A.~J. (1996).
Satisfaction and comparison income.
\textit{Journal of Public Economics}, 61(3), 359--381.

\bibitem[Dahl and Ransom(1999)]{dahl1999}
Dahl, G.~B., \& Ransom, M.~R. (1999).
Does where you stand depend on where you sit? Tithing donations and self-serving beliefs.
\textit{American Economic Review}, 89(4), 703--727.

\bibitem[Fliessbach et al.(2007)]{fliessbach2007}
Fliessbach, K., Weber, B., Trautner, P., Dohmen, T., Sunde, U., Elger, C.~E., \& Falk, A. (2007).
Social comparison affects reward-related brain activity in the human ventral striatum.
\textit{Science}, 318(5854), 1305--1308.

\bibitem[Fong(2001)]{fong2001}
Fong, C. (2001).
Social preferences, self-interest, and the demand for redistribution.
\textit{Journal of Public Economics}, 82(2), 225--246.

\bibitem[Haisley and Weber(2010)]{haisley2010}
Haisley, E.~C., \& Weber, R.~A. (2010).
Self-serving interpretations of ambiguity in other-regarding behavior.
\textit{Games and Economic Behavior}, 68(2), 614--625.

\bibitem[Imai et al.(2010)]{imai2010}
Imai, K., Keele, L., \& Tingley, D. (2010).
A general approach to causal mediation analysis.
\textit{Psychological Methods}, 15(4), 309--334.

\bibitem[Konow(2000)]{konow2000}
Konow, J. (2000).
Fair shares: Accountability and cognitive dissonance in allocation decisions.
\textit{American Economic Review}, 90(4), 1072--1091.

\bibitem[Kunda(1990)]{kunda1990}
Kunda, Z. (1990).
The case for motivated reasoning.
\textit{Psychological Bulletin}, 108(3), 480--498.

\bibitem[Meltzer and Richard(1981)]{meltzer1981}
Meltzer, A.~H., \& Richard, S.~F. (1981).
A rational theory of the size of government.
\textit{Journal of Political Economy}, 89(5), 914--927.

\bibitem[Piketty(1995)]{piketty1995}
Piketty, T. (1995).
Social mobility and redistributive politics.
\textit{Quarterly Journal of Economics}, 110(3), 551--584.

\end{thebibliography}

\end{document}
